\chapter{Abbreviations and Glossary}
\label{ch:glossary}

\section{Abbreviations}
The abbreviations and acronyms used throughout this thesis are collected in
Table~\ref{tab:abbr}, listed alphabetically with their full forms. They span the
air-quality, statistical, machine-learning and software concepts on which the
project draws.

\begin{xltabular}{\textwidth}{B Y}
\caption{List of abbreviations.}\label{tab:abbr}\\
\toprule
\bfseries Abbreviation & \bfseries Meaning \\
\midrule
\endfirsthead
\toprule
\bfseries Abbreviation & \bfseries Meaning \\
\midrule
\endhead
\bottomrule
\endlastfoot
AJK & Azad Jammu \& Kashmir \\
API & Application Programming Interface \\
AQI & Air Quality Index \\
ARIMA & Autoregressive Integrated Moving Average \\
CAMS & Copernicus Atmosphere Monitoring Service \\
CI/CD & Continuous Integration and Continuous Deployment \\
CO & Carbon Monoxide \\
CORS & Cross-Origin Resource Sharing \\
CPU & Central Processing Unit \\
CSS & Cascading Style Sheets \\
EPA & (United States) Environmental Protection Agency \\
ERA5 & The fifth-generation ECMWF atmospheric reanalysis \\
FR & Functional Requirement \\
FTI & Feature-Training-Inference (the pipeline architecture pattern) \\
GB & Gilgit-Baltistan \\
GPU & Graphics Processing Unit \\
HTTP(S) & HyperText Transfer Protocol (Secure) \\
HUDI & Hadoop Upserts, Deletes and Incrementals (a time-travel table format) \\
ICT & Islamabad Capital Territory \\
JSON & JavaScript Object Notation \\
LLM & Large Language Model \\
LSTM & Long Short-Term Memory (a recurrent neural network architecture) \\
MAE & Mean Absolute Error \\
MCP & Model Context Protocol \\
ML & Machine Learning \\
MLOps & Machine-Learning Operations \\
NFR & Non-Functional Requirement \\
NO\textsubscript{2} & Nitrogen Dioxide \\
O\textsubscript{3} & Ozone \\
OpenMP & Open Multi-Processing (the shared-memory parallelism runtime) \\
PM\textsubscript{2.5} & Particulate Matter of diameter 2.5 micrometres or less \\
PM\textsubscript{10} & Particulate Matter of diameter 10 micrometres or less \\
PSI & Population Stability Index \\
$R^2$ & Coefficient of Determination \\
RMSE & Root Mean Squared Error \\
SDLC & Software Development Life Cycle \\
SHAP & SHapley Additive exPlanations \\
SO\textsubscript{2} & Sulphur Dioxide \\
SPA & Single-Page Application \\
TC & Test Case \\
UC & Use Case \\
UI & User Interface \\
URL & Uniform Resource Locator \\
UTC & Coordinated Universal Time \\
WHO & World Health Organization \\
XGBoost & Extreme Gradient Boosting \\
\end{xltabular}

\section{Glossary}
\begin{description}[style=nextline,leftmargin=1.5em,labelsep=0pt,itemsep=5pt]
  \item[Air Quality Index] A dimensionless public-health index computed from
  pollutant concentrations by piecewise-linear interpolation, reported as the
  maximum of the per-pollutant sub-indices and divided into six named health
  categories.
  \item[Anchor state] The observed pollution conditions at the moment a forecast
  is issued - the current index, its recent means and volatility, and the current
  particulate concentrations - supplied to the model as the state from which the
  forecast departs.
  \item[Backfill] The one-off pipeline run that loads the full multi-year history
  into the feature store, in chunks, resumably, so that a model can be trained on
  more than the few days the hourly pipeline has accumulated.
  \item[Baseline (persistence)] The naive forecast that the value $h$ hours ahead
  equals the value now. Scored in every training run as the reference against
  which any claim of forecasting skill must be made.
  \item[Challenger] The best fitted candidate produced by a training run, which
  becomes the champion only if it improves on the incumbent's validation error.
  \item[Champion] The model currently in service. Defined operationally as the
  registered version with the lowest recorded validation RMSE.
  \item[Champion-challenger gate] The rule that a newly trained model may replace
  the model in service only on a measured improvement, so that an unattended
  retrain cannot degrade the system.
  \item[Chronological split] Dividing training from validation by time rather than
  at random, so that the model is always evaluated on a period it has not seen.
  \item[Conformal prediction] A distribution-free framework for attaching
  prediction intervals to an arbitrary regressor; in its split form, the interval
  is built from the empirical quantiles of residuals on a held-out calibration
  set.
  \item[Cyclical encoding] Representing a periodic quantity such as the hour of
  the day by its sine and cosine, so that the last value of the cycle is adjacent
  to the first rather than maximally distant from it.
  \item[Direct multi-horizon forecasting] Predicting the value at $t+h$ directly,
  rather than applying a one-step model repeatedly to its own output, thereby
  avoiding the accumulation of error across steps.
  \item[Drift] A change over time in the distribution of a model's inputs, or in
  the relationship between inputs and target. For a seasonal system, recurring
  drift is the normal condition rather than a fault.
  \item[Ephemeral runner] A scheduled compute instance that is created for one job
  and destroyed when it ends, so that nothing written to its local disk survives.
  \item[Event time] The column recording when an observation actually occurred, as
  distinct from when it was written; used by the feature store to order and
  time-travel over rows.
  \item[Feature group] The versioned, schema-defined table in the feature store
  into which the feature pipeline writes and from which the training pipeline
  reads.
  \item[Feature store] A managed repository of engineered features with primary
  keys and event time, written by feature pipelines and read by both training and
  inference, which eliminates training-serving skew by construction.
  \item[Feature-Training-Inference (FTI)] The architectural pattern in which a
  machine-learning system is expressed as three pipelines that share no code path
  and communicate only through a feature store and a model registry.
  \item[Global model] A single model trained across all locations and all
  horizons, conditioned on location and horizon as features, rather than one model
  per location or per horizon.
  \item[Horizon] The forecast lead time in hours. In this system it is an input
  feature, which is what allows one model to serve every lead time.
  \item[Idempotent upsert] A write that inserts a row if its key is new and
  replaces it if the key already exists, so that repeating the write has no
  additional effect and re-running a pipeline is always safe.
  \item[Leaderboard] The persistent record of every training run, its candidates,
  its metrics and its promotion decision, serving as the audit trail of the
  system's model history.
  \item[Model registry] The durable, versioned store of trained model artefacts
  tagged with their metrics, from which the inference pipeline retrieves the
  champion; the mechanism by which a model outlives the machine that trained it.
  \item[Population Stability Index] A non-parametric measure of how far a current
  distribution has moved from a reference distribution, computed over shared bins;
  conventionally read as stable below 0.10, moderate to 0.25, and significant
  above.
  \item[Prediction interval] A range around a point forecast expected to contain
  the true value with a stated probability; here the tenth to ninetieth residual
  percentiles for the relevant horizon, giving a nominal 80\% interval.
  \item[Serverless] An architecture in which no long-running server is
  provisioned; compute is scheduled, short-lived and disposable, and all durable
  state lives in managed services.
  \item[Shapley value] The unique attribution of a prediction among its input
  features that satisfies local accuracy, missingness and consistency; computable
  exactly and in polynomial time for tree ensembles.
  \item[Target leakage] The presence in a feature of information that would not
  have been available at prediction time, which inflates measured accuracy and
  destroys real-world performance.
  \item[Target-time features] The inputs describing the hour being forecast -
  forecast weather, deterministic calendar and fixed location - as distinct from
  the anchor state describing the moment the forecast is issued.
  \item[Training-serving skew] The failure mode in which features are computed
  differently for training and for inference, so that a model performs worse in
  production than in evaluation for reasons unrelated to the data.
  \item[Walk-forward backtest] Rolling validation in which the training window
  expands and the test window moves forward with it, so that every evaluation is
  made on the future of its own training data.
  \item[What-if simulation] Re-predicting under deliberately altered inputs to
  observe how the model responds; a statement about the model, not a causal claim
  about the world.
\end{description}
